\documentclass[12pt]{article}
\usepackage[spanish]{babel}
\usepackage[utf8]{inputenc}
\usepackage{amsmath, amssymb, amsthm}
\usepackage[margin=2cm]{geometry}
\usepackage{booktabs}
\usepackage{enumitem}
\usepackage{tikz}
\usepackage{float}
\usepackage[hidelinks]{hyperref}

\newcommand{\inputfigura}[1]{\IfFileExists{Parcial2/Figuras/#1.tex}{\input{Parcial2/Figuras/#1.tex}}{\input{Figuras/#1.tex}}}

\begin{document}

\begin{center}
    {\Large \textbf{UNIVERSIDAD DE ANTIOQUIA}}\\[0.2cm]
    {\large \textbf{Astroestadística}}\\[0.1cm]
    \textbf{Problemas para Evaluación \#2}
\end{center}

\vspace{0.5cm}

\section*{Punto 1: Conceptos de estimación e hipótesis}

\subsection*{Enunciado}
Responda las siguientes preguntas:
\begin{enumerate}[label=(\alph*)]
    \item ¿Cuáles son las características de un buen estimador? Descríbalas.
    \item ¿Cómo se define la función de Likelihood? Suponga una muestra independiente e idénticamente distribuida.
    \item ¿En qué consiste la estimación por momentos?
    \item En una prueba de hipótesis, ¿qué es el valor-$p$ y cómo se relaciona con la regla de decisión?
\end{enumerate}

\subsection*{Solución}
\begin{enumerate}[label=(\alph*)]
    \item Un buen estimador busca acercarse al parámetro verdadero y variar poco entre muestras. Algunas propiedades importantes son:
    \begin{itemize}
        \item \textbf{Insesgadez:} $E(\hat{\theta})=\theta$; en promedio, el estimador coincide con el valor del parámetro.
        \item \textbf{Consistencia:} $\hat{\theta}$ se acerca a $\theta$ al aumentar el tamaño muestral; formalmente, para todo $\varepsilon>0$, $P(|\hat{\theta}-\theta|>\varepsilon)\to0$ cuando $n\to\infty$.
        \item \textbf{Eficiencia:} entre estimadores comparables, suele preferirse el que tiene menor varianza. Por ejemplo, al comparar estimadores insesgados, menor $\operatorname{Var}(\hat{\theta})$ significa mayor precisión.
    \end{itemize}

    \item Dada una muestra i.i.d. $x_1,x_2,\ldots,x_n$ con función de densidad o masa $f(x\mid\theta)$, la función de verosimilitud es la función del parámetro $\theta$ que se obtiene al fijar los datos observados:
    $$
    L(x_1,x_2,\ldots,x_n\mid\theta)
    =\prod_{i=1}^{n}f(x_i\mid\theta)
    =f(x_1\mid\theta)f(x_2\mid\theta)\cdots f(x_n\mid\theta).
    $$
    El producto se debe a la independencia de las observaciones. La verosimilitud muestra qué valores de $\theta$ describen mejor los datos, pero no es por sí misma una distribución de probabilidad de $\theta$.

    \item La estimación por momentos consiste en igualar momentos poblacionales con los momentos muestrales correspondientes y resolver las ecuaciones para obtener los parámetros desconocidos. Para el momento de orden $k$, se igualan
    $$
    E(X^k)=\frac{1}{n}\sum_{i=1}^{n}x_i^k.
    $$
    Si hay varios parámetros desconocidos, se utilizan tantos momentos como ecuaciones sean necesarias para determinarlos.

    \item El valor-$p$ es la probabilidad, suponiendo cierta $H_0$, de obtener un resultado tan extremo o más extremo que el observado, de acuerdo con la hipótesis alternativa. No es la probabilidad de que $H_0$ sea cierta.

    Para una prueba de nivel $\alpha$, se rechaza $H_0$ si el valor-$p$ es menor o igual que $\alpha$; en caso contrario, no se rechaza. Si $F$ es la función de distribución acumulada bajo $H_0$, entonces:
    \begin{itemize}
        \item Cola derecha ($\theta>\theta_0$): $\text{Valor-}p=1-F(Z_{\mathrm{obs}})$.
        \item Cola izquierda ($\theta<\theta_0$): $\text{Valor-}p=F(Z_{\mathrm{obs}})$.
        \item Dos colas: si los límites observados de la región extrema son $Z_{\mathrm{obs}}^{-}$ y $Z_{\mathrm{obs}}^{+}$, el valor-$p$ es
        $$
        F(Z_{\mathrm{obs}}^{-})+1-F(Z_{\mathrm{obs}}^{+}).
        $$
        Para una distribución simétrica y una prueba bilateral usual, esto se reduce a $2[1-F(|Z_{\mathrm{obs}}|)]$.
    \end{itemize}
\end{enumerate}

\section*{Punto 2: Biofirmas}

\subsection*{Enunciado}
Investigadores de la UdeA estudian la presencia de una molécula atmosférica $S$ en exoplanetas rocosos, que puede asociarse a procesos biológicos como la fotosíntesis. Sin embargo, procesos abióticos también pueden producirla. Se comparan dos muestras independientes:
\begin{itemize}
    \item \textbf{Planetas dentro de la zona de habitabilidad (HZ):} $n_1=16$, $\bar{x}_1=2.85\,\mathrm{ppm}$ y $s_1=0.50\,\mathrm{ppm}$.
    \item \textbf{Planetas fuera de la zona de habitabilidad:} $n_2=20$, $\bar{x}_2=2.35\,\mathrm{ppm}$ y $s_2=0.45\,\mathrm{ppm}$.
\end{itemize}
Responda:
\begin{enumerate}[label=(\alph*)]
    \item En el contexto del problema, ¿qué representa un error de tipo I y qué representa uno de tipo II?
    \item Realice una prueba de hipótesis para determinar si los planetas en HZ presentan, en promedio, concentraciones mayores de $S$ que los planetas fuera de ella. Utilice $\alpha=0.05$.
    \item Suponga que la diferencia verdadera entre las concentraciones promedio es $\mu_1-\mu_2=0.30\,\mathrm{ppm}$ y que la desviación estándar poblacional común es $\sigma=0.48\,\mathrm{ppm}$. Calcule la probabilidad de error de tipo II ($\beta$) y la potencia de la prueba.
\end{enumerate}

\subsection*{Solución}
\begin{enumerate}[label=(\alph*)]
    \item Un \textbf{error de tipo I} consiste en concluir que la concentración promedio en HZ es mayor cuando en realidad no lo es. Un \textbf{error de tipo II} consiste en no detectar que la concentración promedio en HZ es mayor cuando en realidad sí lo es. Estos errores describen decisiones sobre los promedios y no prueban ni descartan por sí solos la existencia de vida.

    \item Definimos las variables y los parámetros:
    \[
    \begin{aligned}
    X_1&:\text{ concentración de }S\text{ en planetas HZ},\\
    X_2&:\text{ concentración de }S\text{ en planetas fuera de HZ}.
    \end{aligned}
    \]
    $$
    \mu_1=E(X_1),\qquad \mu_2=E(X_2),\qquad
    \Delta=\mu_1-\mu_2.
    $$
    La pregunta es si $\mu_1>\mu_2$. Planteamos las hipótesis
    $$
    H_0:\mu_1-\mu_2=0=\Delta_0,
    \qquad
    H_a:\mu_1-\mu_2>0.
    $$
    La hipótesis nula compuesta incluye los casos $\mu_1-\mu_2\leq0$; para construir la prueba unilateral evaluamos su frontera, $\Delta_0=0$.

    Usamos el estadístico $T$ con los datos de las dos muestras. Aquí $S_1=0.50\,\mathrm{ppm}$ y $S_2=0.45\,\mathrm{ppm}$ son las desviaciones estándar muestrales, y $\hat{\mu}_1=\bar{x}_1$, $\hat{\mu}_2=\bar{x}_2$. El error estándar estimado de la diferencia es
    \begin{align*}
    SE
    &=\sqrt{\frac{S_1^2}{n_1}+\frac{S_2^2}{n_2}}\\
    &=\sqrt{\frac{0.50^2}{16}+\frac{0.45^2}{20}}\\
    &=\sqrt{\frac{0.25}{16}+\frac{0.2025}{20}}\\
    &=\sqrt{0.015625+0.010125}\\
    &=\sqrt{0.02575}\\
    &\approx0.16047\,\mathrm{ppm}.
    \end{align*}

    El estadístico de prueba es
    $$
    T=\frac{\hat{\mu}_1-\hat{\mu}_2-\Delta_0}
    {\sqrt{\dfrac{S_1^2}{n_1}+\dfrac{S_2^2}{n_2}}}
    =\frac{\bar{x}_1-\bar{x}_2-\Delta_0}{SE}.
    $$
    Sustituyendo las medias y el error estándar:
    \begin{align*}
    t_{\mathrm{obs}}
    &=\frac{2.85-2.35-0}{0.16047}\\
    &=\frac{0.50}{0.16047}\\
    &\approx3.116.
    \end{align*}

    Usamos la aproximación solicitada para los grados de libertad:
    $$
    \nu\approx n_1+n_2-2=16+20-2=34.
    $$
    Para una prueba de cola derecha con $\alpha=0.05$, el valor crítico satisface $P(t_{34}>t_{0.05,34})=0.05$. De la distribución $t$:
    $$
    t_{0.05,34}\approx1.691.
    $$
    Como $t_{\mathrm{obs}}=3.116>1.691$, rechazamos $H_0$. El valor-$p$ unilateral es
    $$
    p=P(t_{34}\geq3.116)\approx0.00186<0.05.
    $$
    Ambos criterios conducen a la misma decisión: hay evidencia estadística de que la concentración promedio de $S$ es mayor en los planetas HZ.

    \inputfigura{prueba_t}

    \item Para calcular $\beta$ utilizamos la información nueva del enunciado: la desviación estándar poblacional común es conocida, $\sigma=0.48\,\mathrm{ppm}$. La diferencia verdadera indicada es
    $$
    \mu_1-\mu_2=\Delta'=0.30\,\mathrm{ppm}.
    $$
    La desviación estándar del estimador de la diferencia de medias es
    \begin{align*}
    \sigma_{\hat{\Delta}}
    &=\sigma_{\bar{x}_1-\bar{x}_2}\\
    &=\sqrt{\frac{\sigma_1^2}{n_1}+\frac{\sigma_2^2}{n_2}}\\
    &=\sqrt{\frac{0.48^2}{16}+\frac{0.48^2}{20}}\\
    &=0.48\sqrt{\frac{1}{16}+\frac{1}{20}}\\
    &\approx0.16099\,\mathrm{ppm}.
    \end{align*}

    Para una prueba unilateral de nivel $\alpha=0.05$, el valor crítico normal es $Z_{\alpha}=1.645$. La región de rechazo puede escribirse en la escala de $\hat{\Delta}$ como $\hat{\Delta}>c$, donde
    \begin{align*}
    c
    &=\Delta_0+Z_{\alpha}\sigma_{\hat{\Delta}}\\
    &=0+1.645(0.16099)\\
    &\approx0.2648\,\mathrm{ppm}.
    \end{align*}
    Por tanto, el error de tipo II ocurre si no se rechaza $H_0$, es decir, si $\hat{\Delta}\leq c$ cuando la diferencia verdadera es $\Delta'$:
    \begin{align*}
    \beta(\Delta')
    &=P(\hat{\Delta}\leq c\mid\Delta=\Delta')\\
    &=P\left(\frac{\hat{\Delta}-\Delta_0}{\sigma_{\hat{\Delta}}}
      \leq Z_{\alpha}\;\middle|\;\Delta=\Delta'\right)\\
    &=P\left(\hat{\Delta}\leq Z_{\alpha}\sigma_{\hat{\Delta}}+\Delta_0
      \;\middle|\;\Delta=\Delta'\right)\\
    &=P\left(\hat{\Delta}-\Delta'
      \leq Z_{\alpha}\sigma_{\hat{\Delta}}+\Delta_0-\Delta'
      \;\middle|\;\Delta=\Delta'\right)\\
    &=P\left(\frac{\hat{\Delta}-\Delta'}{\sigma_{\hat{\Delta}}}
      \leq Z_{\alpha}+\frac{\Delta_0-\Delta'}{\sigma_{\hat{\Delta}}}
      \;\middle|\;\Delta=\Delta'\right)\\
    &=\Phi\left(Z_{\alpha}+\frac{\Delta_0-\Delta'}{\sigma_{\hat{\Delta}}}\right)\\
    &=\Phi\left(1.645-\frac{0.30}{0.16099}\right)\\
    &\approx\Phi(-0.2186)\\
    &\approx0.4135.
    \end{align*}
    En consecuencia, la potencia de la prueba es
    $$
    1-\beta=1-0.4135=0.5865.
    $$
    Bajo la diferencia supuesta de $0.30\,\mathrm{ppm}$, la prueba tiene aproximadamente $41.35\%$ de probabilidad de no rechazar $H_0$ y $58.65\%$ de probabilidad de rechazarla.

    \inputfigura{beta_potencia}
\end{enumerate}

\section*{Punto 3: Binarias eclipsantes}

\subsection*{Enunciado}
Algunas bajas en la curva de luz de una estrella pueden atribuirse a un tránsito planetario. Sin embargo, una estrella compañera que pasa por delante o por detrás de la otra también puede producir estos cambios; este fenómeno se conoce como binaria eclipsante. En una muestra de $88$ curvas reportadas como tránsitos, se encontró que $31$ eran binarias eclipsantes.
\begin{enumerate}[label=(\alph*)]
    \item Calcule e interprete un intervalo de confianza (IC) del $95\%$ para la proporción real de binarias eclipsantes entre todas las curvas reportadas como tránsitos bajo estas circunstancias.
    \item ¿Qué tamaño de muestra se requiere para que el ancho deseado del IC aproximado del $95\%$ sea como máximo $0.04$, independientemente de los resultados de la muestra?
\end{enumerate}

\subsection*{Solución}
\begin{enumerate}[label=(\alph*)]
    \item Sea $p$ la proporción poblacional de curvas que corresponden a binarias eclipsantes. En la muestra, $x=31$ de las $n=88$ curvas son binarias eclipsantes, por lo que
    $$
    \hat{p}=\frac{x}{n}=\frac{31}{88}\approx0.35227.
    $$
    Calculamos un intervalo de confianza de Wilson del $95\%$. Con $z_{0.975}=1.96$, los límites se obtienen mediante
    $$
    \frac{\hat{p}+z^2/(2n)}{1+z^2/n}
    \;\pm\;
    \frac{z}{1+z^2/n}
    \sqrt{\frac{\hat{p}(1-\hat{p})}{n}+\frac{z^2}{4n^2}}.
    $$
    Sustituimos $\hat{p}=31/88\approx0.35227$, $n=88$ y $z=1.96$. El centro del intervalo es
    $$
    \frac{0.35227+1.96^2/(2\cdot88)}{1+1.96^2/88}
    \approx0.35845,
    $$
    y el semiancho es
    $$
    \frac{1.96}{1+1.96^2/88}
    \sqrt{\frac{0.35227(1-0.35227)}{88}+\frac{1.96^2}{4\cdot88^2}}
    \approx0.09789.
    $$
    Así, el intervalo de confianza es
    $$
    \boxed{IC_{95\%}(p)\approx[0.26056,\,0.45634]}.
    $$
    En porcentaje, va aproximadamente de $26.06\%$ a $45.63\%$. Su interpretación es que el procedimiento usado produce intervalos que contienen la proporción poblacional en aproximadamente el $95\%$ de las muestras repetidas; el $95\%$ no representa una probabilidad asignada a $p$ después de observar este intervalo.

    \inputfigura{intervalo_binomial}

    \item El ancho máximo deseado es $0.04$, de modo que el margen de error debe satisfacer $E=0.04/2=0.02$. Para la aproximación normal al $95\%$,
    $$
    E=z_{0.975}\sqrt{\frac{p(1-p)}{n}}.
    $$
    Como $p$ es desconocido y el tamaño de muestra no debe depender de los resultados, usamos el peor caso $p(1-p)\leq1/4$, que ocurre en $p=0.5$. Entonces,
    \begin{align*}
    1.96\sqrt{\frac{0.25}{n}}&\leq0.02,\\
    \frac{1.96^2(0.25)}{n}&\leq0.02^2,\\
    n&\geq\frac{1.96^2(0.25)}{0.02^2}=2401.
    \end{align*}
    Redondeando hacia arriba al entero, el tamaño mínimo requerido es
    $$
    \boxed{n_{\min}=2401.}
    $$
    Este resultado usa la aproximación normal y el caso más conservador $p=0.5$; el ancho obtenido con otro método de intervalo puede variar.

    \inputfigura{tamano_muestra}
\end{enumerate}

\end{document}